%===================================================================
% Chapter 5: Future Work -- Phase-1 stub. Seeded from directions
% mentioned across the three papers; expand and revise.
%===================================================================
\chapter{Future Work}
\label{chap:future}

% TODO: expand each direction into a full paragraph or section.
This dissertation opens several directions for future research.

\paragraph{From rewrite discovery to rule discovery.}
The dataflow-guided synthesis framework of SlabCity could be applied
to discovering reusable rewrite \emph{rules}, complementing systems
such as WeTune~\cite{wetune}. Similarly, NLR2s distilled by GenRewrite
form a growing, human-readable rule repository; studying how such
repositories transfer across schemas, engines, and organizations is a
natural next step.

\paragraph{Tighter integration with query optimizers.}
% TODO: flesh out -- e.g., feeding rewriter signals into cost models,
% or invoking higher-level rewriters selectively based on optimizer
% confidence.
Our systems currently operate as source-to-source preprocessors
outside the database. Integrating them more tightly with the
optimizer could reduce overhead and enable cooperative optimization.

\paragraph{Beyond performance: cost, correctness, and maintainability.}
% TODO: flesh out -- pipeline-level objectives beyond latency, e.g.,
% cloud cost budgets, freshness constraints, and developer-friendly
% refactorings, extending the DAGSmith direction.
DAGSmith suggests that pipeline-level rewriting must balance latency,
warehouse cost, refresh frequency, and the readability of the
resulting SQL project; formalizing these trade-offs is an open
problem.
